\begin{tikzpicture}[font=\small]
  % ---------- magnitude response vs order ----------
  \begin{scope}
    \draw[ax] (-0.2,0) -- (5.4,0); \node[note, anchor=north west] at (5.2,-0.05){$\omega/\omega_c$};
    \draw[ax] (0,-0.2) -- (0,2.4); \node[note, anchor=south east] at (-0.1,2.2){$|H(j\omega)|$};

    \draw[helper] (0,1.9) -- (5.2,1.9);
    \node[note, anchor=east] at (-0.1,1.9){$1$};
    \draw[helper] (0,1.343) -- (5.2,1.343);
    \node[note, anchor=east] at (-0.1,1.343){$\tfrac{1}{\sqrt2}$};
    \draw[helper] (2.0,0) -- (2.0,2.0);
    \node[note, anchor=north] at (2.0,-0.1){$1$};

    \foreach \N/\col/\lbl in {1/sslight/{$N=1$}, 2/ssteal/{$N=2$}, 4/ssblue/{$N=4$}, 8/ssred/{$N=8$}}{
      \draw[\col, line width=1.4pt, domain=0:5.2, samples=300, smooth]
        plot (\x, {1.9*\bwmag{\x/2.0}{\N}});
    }
    \node[sslight, font=\footnotesize, anchor=west] at (4.65,0.88){$N=1$};
    \node[ssteal,  font=\footnotesize, anchor=west] at (4.65,0.30){$N=2$};
    \node[ssblue,  font=\footnotesize, anchor=west] at (2.95,0.62){$N=4$};
    \node[ssred,   font=\footnotesize, anchor=east] at (1.35,1.95){$N=8$};

    \node[note, anchor=north, align=center] at (2.6,-0.6)
      {$|H(j\omega)|^2=\dfrac{1}{1+(\omega/\omega_c)^{2N}}$\\[3pt]
       通带内\textbf{最大平坦}（前 $2N-1$ 阶导数在 $\omega=0$ 处为零），\\
       无纹波；代价是过渡带比切比雪夫、椭圆都宽};
  \end{scope}

  % ---------- pole locations ----------
  \begin{scope}[xshift=9.6cm]
    \draw[ax] (-3.2,0) -- (3.0,0); \node[note, anchor=north west] at (2.8,-0.05){$\sigma$};
    \draw[ax] (0,-2.9) -- (0,2.9); \node[note, anchor=south west] at (0.08,2.75){$j\omega$};
    \draw[helper] (0,0) circle (2.2);

    % N = 4 : poles at angles 112.5, 157.5, 202.5, 247.5 degrees
    \foreach \a in {112.5,157.5,202.5,247.5}{
      \draw[ssblue, line width=1.4pt]
        ({2.2*cos(\a)-0.15},{2.2*sin(\a)-0.15}) -- ({2.2*cos(\a)+0.15},{2.2*sin(\a)+0.15});
      \draw[ssblue, line width=1.4pt]
        ({2.2*cos(\a)-0.15},{2.2*sin(\a)+0.15}) -- ({2.2*cos(\a)+0.15},{2.2*sin(\a)-0.15});
    }
    % mirrored (unstable) ones shown faintly
    \foreach \a in {-22.5,22.5,-67.5,67.5}{
      \draw[sslight, line width=1.0pt]
        ({2.2*cos(\a)-0.13},{2.2*sin(\a)-0.13}) -- ({2.2*cos(\a)+0.13},{2.2*sin(\a)+0.13});
      \draw[sslight, line width=1.0pt]
        ({2.2*cos(\a)-0.13},{2.2*sin(\a)+0.13}) -- ({2.2*cos(\a)+0.13},{2.2*sin(\a)-0.13});
    }
    \node[ssblue, font=\footnotesize, anchor=east] at (-1.55,2.35){取左半平面这 $N$ 个};
    \node[note, font=\footnotesize, anchor=west] at (1.35,2.35){右半平面的镜像丢弃};
    \node[note, anchor=north east] at (2.15,-0.15){$\omega_c$};

    \node[note, anchor=north, align=center] at (0,-3.3)
      {$2N$ 个极点\textbf{等角均布}在半径 $\omega_c$ 的圆上；\\
       取左半平面那 $N$ 个即得因果稳定的 $H(s)$。\\
       $N$ 越大极点越密集地贴近 $j\omega$ 轴，滚降越陡但越接近振荡};
  \end{scope}
\end{tikzpicture}
